\chapter{Recommandations pour les TPs}

Le but de ces travaux pratiques est d’illustrer le cours de conversion d’énergie de deuxième année par l’étude de quatre manipulations :

\begin{itemize}
\item Hacheur Buck (simulations avec LTSpice)
\item Hacheur 4 quadrants
\item Alimentation à découpage Flyback
\item Alimentation à découpage Forward
\end{itemize}

Afin de préparer ces séances, il est nécessaire de connaître le cours correspondant et d’avoir lu le texte complet du TP afin de prévoir les différentes étapes de l’étude. En particulier, lorsqu'une comparaison avec les courbes théoriques est demandée, il est nécessaire de les apporter en séance et de savoir les adapter au TP.

Les thèmes de TP se déroulent en rotation : il est donc impératif de connaître son numéro de binôme avant de venir en TP.

\vspace{0.5cm}
\hspace{-1cm}
\begin{tabular}{|c | c | c | c | c | c | c | c|}
\hline
& \textbf{Binôme 1} & \textbf{Binôme 2} & \textbf{Binôme 3} & \textbf{Binôme 4} & \textbf{Binôme 5} & \textbf{Binôme 6} & \textbf{Binôme 7}\\
\hline
\textit{Séance 1} & Hach. 4Q & Hach. 4Q & Hach. 4Q & LTSpice & LTSpice & LTSpice & LTSpice\\
\hline
\textit{Séance 2} & LTSpice & LTSpice & LTSpice & Hach 4Q & Hach. 4Q & Hach. 4Q & Hach. 4Q\\
\hline
\textit{Séance 3} & Forward & Forward & Forward & Flyback & Flyback & Flyback & Flyback\\
\hline
\textit{Séance 4} & Flyback & Flyback & Flyback & Forward & Forward & Forward & Forward\\
\hline
\end{tabular}

\vspace{0.5cm}
Suivant le laboratoire où se déroulent les séances de TP, quelques différences peuvent apparaître sur le câblage des machines lors du TP hacheur quatre quadrants ainsi que sur les grandeurs nominales de ces machines qui figurent sur leur plaque signalétique et qui devront donc être soigneusement notées.

\vspace{0.5cm}
\textbf{\underline{A l’issue de chaque séance}, un compte-rendu devra être rendu au professeur. L’évaluation tiendra compte à la fois de la qualité du compte rendu et du comportement lors de la séance elle-même.}